\subsection{Emergence of correlations}
  \label{subsec:emergence-of-correlations}

  Let $\mu(\omega)$ be a probability distribution defined on $\Omega$.
  The joint probability distribution accessible at the observable level is given by
  \[
    P(a,b) = \sum_{\omega \in \Pi^{-1}(a,b)} \mu(\omega).
  \]

  Because the preimage $\Pi^{-1}(a,b)$ does not generically decompose into independent
  subsets associated with $A$ and $B$, the joint probability $P(a,b)$ does not admit a
  factorized representation of the form $P(a)P(b)$ or
  $P(a \mid \lambda) P(b \mid \lambda)$ for any effective variable $\lambda$.

  The resulting correlations arise not from any interaction or information exchange
  between $A$ and $B$, but from the fact that both outcomes are constrained by the same
  underlying relational configuration.
  From the perspective of the effective description, the two outcomes appear correlated
  even when the corresponding measurements are spacelike separated.
